\noindent\textbf{উদাহরণ:} $(6B.9)_{16}$-কে অক্টালে রূপান্তর করো।

\cvStep{১}{প্রতিটি হেক্সাডেসিমাল অঙ্কের জন্য চারটি বিট}
\begin{center}
\begin{tikzpicture}[x=1mm,y=1mm]
 \cvCanvas{-34}
 \cvDigit{hex-six}{38}{0}{6}
 \cvDigit{hex-b}{70}{0}{B}
 \cvDigit{hex-nine}{118}{0}{9}
 \node[cv digit] at (94,0) {$.$};
 \cvBits{hex-bits}{-18}{26/0,34/1,42/1,50/0,58/1,66/0,74/1,82/1,106/1,114/0,122/0,130/1}
 \node[cv digit] at (94,-18) {$.$};
 \cvMapArrow{hex-six}{38}{-14}
 \cvMapArrow{hex-b}{70}{-14}
 \cvMapArrow{hex-nine}{118}{-14}
 \cvFourBrace{38}
 \cvFourBrace{70}
 \cvFourBrace{118}
\end{tikzpicture}
\end{center}
\noindent $6=0110_2$, $B=11_{10}=1011_2$, $9=1001_2$। বিটের ক্রম ঠিক রেখে:
\[
 (6B.9)_{16}=(01101011.1001)_2.
\]

\cvStep{২}{বিন্দু থেকে দুই দিকে তিন বিট করে নতুন গ্রুপ}
\noindent পূর্ণাংশের বামে একটি ০ এবং ভগ্নাংশের ডানে দুটি ০ বসাও।
\begin{center}
\begin{tikzpicture}[x=1mm,y=1mm]
 \cvCanvas{-45}
 % Unpadded row and padded row share the same bit columns and radix point.
 \cvBits{original}{0}{26/0,34/1,42/1,50/0,58/1,66/0,74/1,82/1,106/1,114/0,122/0,130/1}
 \node[cv digit] at (94,0) {$.$};
 \cvBits{kept}{-18}{26/0,34/1,42/1,50/0,58/1,66/0,74/1,82/1,106/1,114/0,122/0,130/1}
 \foreach \i in {1,...,12}{\draw[cv arrow,draw=black!65] (original-\i.south)--(kept-\i.north);}
 \cvDigit{pad-left}{18}{-18}{0}
 \cvDigit{pad-right-one}{138}{-18}{0}
 \cvDigit{pad-right-two}{146}{-18}{0}
 \node[cv padding,fit=(pad-left)] (left-zero) {};
 \node[cv padding,fit=(pad-right-one)(pad-right-two)] (right-zeros) {};
 \node[cv note] (left-note) at (6,-5) {বামে ১টি ০};
 \draw[cv arrow] (left-note.south)--(left-zero.north west);
 \node[cv note] (right-note) at (145,-5) {ডানে ২টি ০};
 \draw[cv arrow] (right-note.south)--(right-zeros.north);
 \node[cv digit] at (94,-18) {$.$};
 % Each brace covers exactly three columns; arrows end on its own digit.
 \cvOctalGroup{octal-one}{26}{18}{34}{1}
 \cvOctalGroup{octal-five}{50}{42}{58}{5}
 \cvOctalGroup{octal-three}{74}{66}{82}{3}
 \cvOctalGroup{octal-four-a}{114}{106}{122}{4}
 \cvOctalGroup{octal-four-b}{138}{130}{146}{4}
 \node[cv digit] at (94,-40) {$.$};
\end{tikzpicture}
\end{center}
\noindent দাগঘেরা বিটগুলো নতুন যোগ করা শূন্য। মূল বিটগুলো সোজা নিচে নেমেছে; তাদের স্থান বদলায়নি। প্রতিটি brace-এর তিন বিট থেকে নিচের একটি অক্টাল অঙ্ক পাওয়া যায়।

\cvStep{৩}{অঙ্কগুলো একই ক্রমে একত্র করো}
\[
 \boxed{\therefore\ (6B.9)_{16}=(153.44)_8}
\]
\noindent\textbf{মান যাচাই:}
\[
 6\times16+11+\frac9{16}
 =1\times64+5\times8+3+\frac4{8}+\frac4{64}
 =107.5625.
\]
\noindent পূর্ণাংশের বামে বা ভগ্নাংশের শেষে শূন্য বসালে মান বদলায় না।
